\documentclass[11pt]{article}
\usepackage{amsmath,amssymb,bm,siunitx}
\usepackage[margin=1in]{geometry}
\usepackage{booktabs,longtable,array}
\usepackage[dvipsnames]{xcolor}
\usepackage[colorlinks=true,linkcolor=NavyBlue,citecolor=NavyBlue,urlcolor=NavyBlue]{hyperref}

\newcommand{\ket}[1]{\lvert #1\rangle}
\newcommand{\bra}[1]{\langle #1\rvert}
\newcommand{\avg}[1]{\left\langle #1\right\rangle}
\newcommand{\Pt}{P_2}
\newcommand{\status}[1]{\par\noindent\textbf{\textsf{Status:}} \emph{#1}\par}

\title{Spin-1 Population Dynamics in Polarized Deuteron Targets:\\
Mechanisms, Rate Equations, and Open Questions}
\author{}
\date{}

\begin{document}
\maketitle

\begin{abstract}
These notes collect the physics that governs the level populations of an ensemble of spin-1
nuclei (deuterons) in a dynamically polarized solid. They cover the Hamiltonian and powder
spectrum, the reduction from the density matrix to rate equations, and every population-changing
process relevant to polarized-target work:
\begin{itemize}
\item radio-frequency (RF) saturation and adiabatic fast passage (AFP);
\item spin-lattice relaxation with separate vector and tensor channels;
\item spectral spin diffusion with its three distinct spin-1 flip-flop channels;
\item the nuclear dipolar reservoir (Provotorov theory);
\item dynamic nuclear polarization by the solid effect, cross effect, thermal mixing and the
differential solid effect.
\end{itemize}
The terms are assembled into one orientation- and site-resolved master equation, and we state
its conservation laws and stationary states. Several general consequences for manipulating the
tensor polarization follow directly from the structure of the equations. Throughout, we
distinguish what is established theory, what is an effective closure, and what is not known
for deuterated butanol or ammonia. Numbers quoted for d-butanol come from fits to UVA target
data and are labelled as such.
\end{abstract}

\tableofcontents

%==========================================================================================
\section{Scope, levels of description, and notation}
\label{sec:scope}

The system is a frozen, amorphous solid (e.g.\ deuterated butanol C$_4$D$_9$OD or ND$_3$) at
$T\simeq 1$\,K in $B_0\simeq 2.5$--$5$\,T, doped with paramagnetic centres (chemical radicals
such as TEMPO or EHBA-Cr(V), or irradiation-induced radicals in ammonia) at concentration
$c_e\sim 10^{19}$--$10^{20}$\,cm$^{-3}$ \cite{CrabbMeyer1997,Goertz2004,Meyer2004}. The
deuteron is a spin $I=1$ nucleus with gyromagnetic ratio
$\gamma_d/2\pi = 6.536$\,MHz/T, so $\omega_d/2\pi\simeq 32.7$\,MHz at 5\,T. Its electric
quadrupole moment couples to the electric-field gradient (EFG) of the chemical bond
\cite{CohenReif1957,Abragam1961,Wasylishen2012}.

\paragraph{Three levels of description.}
\begin{enumerate}
\item \emph{Density matrix} of the coupled electron--nuclear system, with relaxation
superoperators (Redfield theory) \cite{Abragam1961,Slichter1990}. This level is exact but
intractable for $\sim 10^{22}$ spins.
\item \emph{Spin-temperature (thermodynamic) theory}: weakly coupled reservoirs (electron
Zeeman, electron dipolar, nuclear Zeeman, nuclear dipolar/quadrupolar), each internally
equilibrated at its own inverse temperature
\cite{Provotorov1962,Goldman1970,AbragamGoldman1978,AbragamGoldman1982}.
\item \emph{Population (rate) equations} for the diagonal elements of the density matrix of
spectrally resolved nuclear ``packets''. Coherences are adiabatically eliminated, and every
process becomes a transition rate or a pair-exchange current. These notes work at this level
and state explicitly where it fails.
\end{enumerate}

\paragraph{Notation.}
\begin{itemize}
\item Levels are $m\in\{+1,0,-1\}$. For a deuteron ($\gamma_d>0$) with
$\mathcal H_Z=-\hbar\omega_d I_z$, the level $m=+1$ is lowest.
\item Populations of the spins at orientation $\Omega$ (EFG principal axes relative to
$\mathbf B_0$, polar angle $\theta$, azimuth $\phi$) and chemical site $s$ are
$n_m(\Omega,s,t)$.
\item Fractional populations within a packet are $p_m$, with $p_{+1}+p_0+p_{-1}=1$.
\item Powder averages are written
$\avg{f}=\frac{1}{4\pi}\int f\,d\Omega$; if $f$ is independent of $\phi$ this reduces to
$\tfrac12\int_0^\pi f\sin\theta\,d\theta$.
\end{itemize}

\paragraph{Polarizations.}
\begin{equation}
P = \avg{n_{+1}-n_{-1}}, \qquad Q=\avg{n_{+1}+n_{-1}-2n_0}=1-3\avg{n_0}.
\label{eq:PQ}
\end{equation}
Positivity of the populations restricts $(P,Q)$ to the triangle with vertices $(\pm1,1)$ and
$(0,-2)$. A Boltzmann distribution with Zeeman inverse spin temperature $\beta$ gives
$n_m\propto e^{m x}$ with $x=\beta\hbar\omega_d$, so that
\begin{equation}
P=\frac{2\sinh x}{1+2\cosh x},\qquad Q=\frac{2\cosh x-2}{1+2\cosh x}=2-\sqrt{4-3P^2}.
\label{eq:boltzmann}
\end{equation}
Corrections of order $\omega_Q/\omega_d$ are discussed in Sec.~\ref{sec:equilibrium}. At
$P=0.50$, Eq.~\eqref{eq:boltzmann} gives $Q=0.1972$ and
$(p_{+1},p_0,p_{-1})=(0.616,0.268,0.116)$; these values are used for illustrations
throughout.

%==========================================================================================
\section{Hamiltonian, transition frequencies and the powder spectrum}
\label{sec:hamiltonian}

\subsection{Zeeman and quadrupolar interactions}
In its principal-axis system the quadrupolar Hamiltonian is
\cite{CohenReif1957,Abragam1961,Wasylishen2012}
\begin{equation}
\mathcal H_Q=\frac{e^2qQ}{4I(2I-1)}\Big[3I_{z'}^2-I(I+1)+\eta\,(I_{x'}^2-I_{y'}^2)\Big],
\end{equation}
where $e^2qQ/h$ is the quadrupolar coupling constant and $0\le\eta\le1$ the EFG asymmetry.
For aliphatic C--D bonds $e^2qQ/h\approx 160$--$170$\,kHz with $\eta\lesssim0.05$; for
hydrogen-bonded O--D, typically 200--230\,kHz \cite{Wasylishen2012}.

In high field, $|\mathcal H_Q|\ll\hbar\omega_d$. To first order only the part commuting with
$I_z$ survives:
\begin{equation}
\mathcal H^{(1)}_Q=\frac{e^2qQ}{4}\,\Pt^\eta(\Omega)\,\big(3I_z^2-2\big)\qquad (I=1),
\qquad
\Pt^\eta(\Omega)\equiv \tfrac12\big(3\cos^2\theta-1+\eta\sin^2\theta\cos2\phi\big).
\end{equation}
The first-order energies are
$E_m=-\hbar\omega_d m+\tfrac{e^2qQ}{4}\Pt^\eta(3m^2-2)$, and the two allowed ($\Delta m=\pm1$)
transitions are
\begin{equation}
\boxed{\;
\omega_{+}(\Omega)\equiv\omega_{+1\leftrightarrow0}=\omega_d-\omega_Q\,\Pt^\eta(\Omega),\qquad
\omega_{-}(\Omega)\equiv\omega_{0\leftrightarrow-1}=\omega_d+\omega_Q\,\Pt^\eta(\Omega),
\qquad \omega_Q\equiv\frac{3e^2qQ}{4\hbar}.\;}
\label{eq:omegapm}
\end{equation}
\noindent\textbf{Convention.} With this definition:
\begin{itemize}
\item the separation of the two transitions at one orientation is $2\omega_Q|\Pt^\eta|$;
\item the "horns" ($\theta=90^\circ$, $\eta=0$) lie at offsets $\pm\omega_Q/2$, and the
outer "shoulders" ($\theta=0$) at $\pm\omega_Q$;
\item both lines coincide at the magic angle ($\cos^2\theta=1/3$).
\end{itemize}
The two forms in the earlier draft differ by a factor of 6 in the definition of $\omega_Q$;
Eq.~\eqref{eq:omegapm} is used everywhere below. For the d-butanol C--D site the fitted horn
offset is $\omega_Q/4\pi=61.7$\,kHz, so $e^2qQ/h=164$\,kHz.

\subsection{Second-order and smaller terms}
\begin{itemize}
\item \textbf{Second-order quadrupole shift.} For $I=1$ it is of order
$\omega_Q^2/\omega_d\sim 0.5$\,kHz at 5\,T for C--D. It shifts \emph{both} transitions by the
same orientation-dependent amount, so it does not change their splitting
\cite{CohenReif1957,Wasylishen2012}. It is below the homogeneous width but not negligible for
O--D.
\item \textbf{Chemical-shift anisotropy} of $^2$H is a few ppm (about 0.1\,kHz): negligible.
\item \textbf{D--D dipolar couplings} are of order 0.1--1\,kHz and set the homogeneous width
(Sec.~\ref{sec:homogeneous}). D--H couplings matter if residual protons are present.
\item \textbf{Hyperfine couplings} to the electrons shift deuterons within $\sim$1\,nm of a
radical far out of the bulk line (Sec.~\ref{sec:barrier}).
\end{itemize}

\subsection{Homogeneous broadening}
\label{sec:homogeneous}
Each packet has a homogeneous width $\Gamma=1/T_2$. It sets three things: the RF response
(Sec.~\ref{sec:rf}), the energy matching of flip-flops (Sec.~\ref{sec:diffusion}), and the local
line $\varphi$ in Eq.~\eqref{eq:line}. Contributions, roughly in order of size for deuterated
glasses at 1\,K:
\begin{enumerate}
\item \textbf{D--D dipolar couplings.} Van Vleck's second moment \cite{VanVleck1948,Abragam1961}
weights neighbours whose transitions are resonant ("like") $9/4$ times more than non-resonant
("unlike") ones. For quadrupole-split deuterons, resonance depends on the neighbour's
orientation, so $\Gamma$ depends on the local density of matching transitions.
\item \textbf{Couplings to residual protons} (heteronuclear, always "unlike").
\item \textbf{Hyperfine fields of the paramagnetic centres.} These give a distribution of shifts,
and they make the width position dependent (largest near radicals).
\item \textbf{Molecular motion} (methyl tunnelling, librations). This can narrow or broaden,
depending on rate.
\item \textbf{Second-order quadrupole dispersion} for sites with a distribution of $\eta$ or
$e^2qQ$. This is inhomogeneous; it is included in the line fit.
\end{enumerate}
The fitted common Lorentzian HWHM for d-butanol, 1.45\,kHz, is an average. Its dependence on
orientation and site is \textbf{not measured}, and it is the main physical candidate for
per-spin RF selectivity (Sec.~\ref{sec:selectivity}).

\subsection{Pake density: the orientation Jacobian}
\label{sec:pake}
For one site with $\eta=0$, write the offset of a transition as
$\delta=\omega-\omega_d=\mp\omega_Q\Pt(\cos\theta)$. The fraction of spins per unit offset is
the orientation Jacobian \cite{Pake1948}:
\begin{equation}
g(\delta)=\frac{1}{|d\delta/d\cos\theta|}=\frac{1}{\sqrt3\,\omega_Q\sqrt{1+2\delta/\omega_Q}},
\qquad \delta\in[-\omega_Q/2,\,\omega_Q],
\label{eq:pakedensity}
\end{equation}
(for the branch whose offset is $+\omega_Q\Pt$; the other branch is mirrored).

\begin{itemize}
\item The density diverges at the horn ($\theta=90^\circ$) and falls to its edge value at
$\theta=0$.
\item With $\eta\neq0$ the three principal values of $\Pt^\eta$ are $1$ and
$-(1\pm\eta)/2$. The divergence moves to $\delta=-(1-\eta)\omega_Q/2$, a new step edge
appears at $-(1+\eta)\omega_Q/2$, and the far edge stays at $\omega_Q$
\cite{CohenReif1957,Wasylishen2012}.
\item Homogeneous broadening convolves each branch with the local line $\varphi$ (Lorentzian
or Voigt). An analytic complex form of the broadened branch is given in
Ref.~\cite{Dulya1997}.
\end{itemize}

Throughout, the \emph{capacity} $\mu$ of a spectral packet is the number of spins whose
transition lies in that packet:
\[
\mu(\delta)\,d\delta = \text{(number of spins)}\times g(\delta)\,d\delta .
\]
\emph{The Jacobian is a density of states. It counts spins; it does not alter the physics of
any single spin.} This distinction matters for the RF rates (Sec.~\ref{sec:rf}).

\subsection{Several sites}
\label{sec:sites}
A molecule such as C$_4$D$_9$OD contains several deuteron sites. Each has its own $e^2qQ$,
$\eta$ and dynamics:
\begin{itemize}
\item CD$_2$ and CD$_3$ groups;
\item one O--D.
\end{itemize}
The spectrum is a weighted sum of Pake patterns. If a methyl group rotates rapidly, its
effective coupling is scaled by $\Pt(\cos109.5^\circ)=-1/3$. At 1\,K, methyl rotation is
frozen or proceeds only by tunnelling, and its effect on the line and on relaxation in
d-butanol is not established.

A fit of 2014 UVA d-butanol spectra with a two-site (C--D, O--D) azimuth-averaged model gave:
\begin{itemize}
\item O--D/C--D coupling ratio 1.35;
\item $\eta_{CD}=0.065$ and $\eta_{OD}=0.033$;
\item an effective O--D integrated weight of 2.1\%;
\item a common Lorentzian HWHM of 1.45\,kHz;
\item held-out RMS of 0.79\% of peak-to-peak.
\end{itemize}
A single-site model gave 0.80\%. The second site is therefore only weakly required, and its
weight should not be read as a chemical abundance.

\subsection{The two overlapping lines}
\label{sec:overlap}
At any frequency in the spectrum two different groups of spins absorb:
\begin{itemize}
\item the I$_+$ transition ($+1\leftrightarrow0$) of spins at orientation $\Omega_1$;
\item the I$_-$ transition ($0\leftrightarrow-1$) of spins at the "mirror" orientation
$\Omega_2$, with $\Pt^\eta(\Omega_2)=-\Pt^\eta(\Omega_1)$.
\end{itemize}
These are different molecules with different capacities. At the negative horn, for example,
the I$_-$ transition belongs to the dense $\theta=90^\circ$ spins, while the overlapping
I$_+$ belongs to sparse spins near $\theta\approx37^\circ$.

\emph{Every manipulation of $Q$ is constrained by this overlap}
(Sec.~\ref{sec:consequences}).

\subsection{Observables}
\label{sec:observables}
The absorptive NMR signal is
\begin{equation}
I(\omega,t)=\sum_s\avg{\big(n_{+1}-n_0\big)\varphi(\omega-\omega_+)+\big(n_0-n_{-1}\big)\varphi(\omega-\omega_-)}_s ,
\label{eq:line}
\end{equation}
with an overall sign set by the detector phase (the earlier draft used the opposite sign
convention). Integrating over each line separately,
\[
A_+=\avg{n_{+1}-n_0},\qquad A_-=\avg{n_0-n_{-1}},\qquad
A_++A_-=P,\qquad A_+-A_-=Q .
\]
\begin{itemize}
\item The total area measures $P$.
\item The area difference of the two branches measures $Q$.
\item In equilibrium the peak-height ratio $r=I_+/I_-=e^{x}$ gives $P$ through
Eq.~\eqref{eq:boltzmann}; this is the basis of the asymmetry method \cite{Dulya1997}.
\end{itemize}
After a manipulation the populations are no longer Boltzmann, so the ratio method fails. In
that case $P$ and $Q$ must be obtained from full line-shape fits.

%==========================================================================================
\section{From the density matrix to rate equations}
\label{sec:reduction}

\subsection{Elimination of coherences}
Take a pair of levels $a,b$ connected by an off-diagonal perturbation $V$ (an RF field, or the
mixing term of a forbidden transition), with detuning $\Delta$ and coherence decay $1/T_2$.
The coherence obeys
\[
\dot c_{ab}=-(1/T_2+i\Delta)c_{ab}+iV(n_a-n_b).
\]
When $|V|T_2\ll1$, or more generally when the coherence decays faster than the populations
change, $c_{ab}$ follows the populations adiabatically. Substituting gives
\cite{Abragam1961,Slichter1990}
\begin{equation}
\dot n_a=-W(n_a-n_b),\qquad
W=\frac{2|V|^2T_2}{1+\Delta^2T_2^2}=2\pi|V|^2 g(\Delta),
\label{eq:golden}
\end{equation}
where $g$ is the normalized homogeneous line. This is Fermi's golden rule with a Lorentzian
line. For a non-Lorentzian homogeneous line, $g$ is replaced by the appropriate
(e.g.\ Gaussian) shape function.

\subsection{Validity}
The rate description requires all of the following:
\begin{enumerate}
\item \textbf{Weak drive}, $\omega_1 T_2\lesssim1$, or at least that population changes be
slow compared with $T_2$.
\item \textbf{Phase-averaged carriers}: no coherent interference between simultaneously
applied tones.
\item \textbf{Secular approximation}: transitions separated by more than their widths are
treated independently.
\end{enumerate}
When condition 1 fails:
\begin{itemize}
\item a resonant pulse produces nutation (at $\sqrt2\,\omega_1$ for a single spin-1
transition) and spin locking;
\item off-resonance irradiation is described by Redfield's rotating-frame thermodynamics
\cite{Redfield1955} or Provotorov's two-temperature theory \cite{Provotorov1962}
(Sec.~\ref{sec:provotorov});
\item quadrupolar echoes and multiple-quantum coherences require the full spin-1 density
matrix \cite{Wasylishen2012}.
\end{itemize}

\subsection{Packet representation}
\label{sec:packets}
Group spins into packets labelled by site $s$ and by the frequency $R$ of their I$_+$
transition. Their I$_-$ transition then lies at the mirror frequency.

\begin{itemize}
\item The state is $n_m(R,s,t)=\mu_s(R)\,p_m(R,s,t)$.
\item All spins in a packet share the same transition frequencies.
\item Packets are therefore the natural variables: orientation enters only through the
capacity $\mu_s(R)$ and through any orientation dependence of the rate coefficients (T$_2$,
couplings).
\end{itemize}

Spins with different azimuth $\phi$ but the same frequency are lumped together when
$\eta\neq0$. This is exact for every process that depends only on frequency. It is an
approximation for processes that depend on the full orientation, such as spatial correlations
of neighbours.

%==========================================================================================
\section{RF saturation (ss-RF) and per-spin selectivity}
\label{sec:rf}

Selective semi-saturating RF (ss-RF) burns spectral holes in the deuteron line to manipulate
the tensor polarization \cite{KellerSSRF,Keller2017}.

\subsection{Weak-drive rates for spin 1}
A linearly polarized RF field contains a rotating component of amplitude $B_1$, with
$\omega_1=\gamma_dB_1$. In the rotating frame $V=\omega_1 I_x$.

\textbf{Matrix elements.} For spin 1,
\[
|\bra{+1}I_x\ket{0}|^2=|\bra{0}I_x\ket{-1}|^2=\tfrac12 .
\]
To first order in $\omega_Q/\omega_d$ these do not depend on orientation, because the
quadrupole interaction does not mix the Zeeman eigenstates at that order.

\textbf{Rate.} From Eq.~\eqref{eq:golden}, a carrier at $\omega_c$ drives a transition at
$\omega_\pm$ with
\begin{equation}
W_\pm=\frac{\omega_1^2T_2}{1+(\omega_\pm-\omega_c)^2T_2^2},\qquad
J_\pm=W_\pm\,\big(n_{\rm lower}-n_{\rm upper}\big).
\label{eq:rfrate}
\end{equation}
\begin{itemize}
\item A carrier exactly on resonance with a two-level pair equalizes it at rate $2W$; with
lattice relaxation the steady state is $d=d_0/(1+2WT_1)$.
\item Independent carriers (random relative phases) add in power: $W=\sum_c W^{(c)}$.
\item RF power is proportional to $B_1^2$, so the rates are linear in power.
\end{itemize}

\subsection{Spin density does not enter the rate per spin}
Equation~\eqref{eq:rfrate} contains $\omega_1$, the detuning and $T_2$. It does \emph{not}
contain how many spins share the frequency $\omega_\pm$:
\begin{itemize}
\item the Pake Jacobian (the horn) puts more spins per hertz near $\theta=90^\circ$;
\item each of those spins sees the same $B_1$ and has the same matrix element;
\item $B_1$ is a field, not a resource shared among the spins: at the RF powers used, spin
absorption is a negligible load on the coil.
\end{itemize}
Hence the absolute rate of polarization loss of a packet is proportional to its capacity $\mu$
(through $n=\mu p$), but its \emph{fractional} saturation rate is not. A phenomenological rate
proportional to the powder density (used in early models) counts the spins twice.

\subsection{Physical sources of per-spin selectivity}
\label{sec:selectivity}
Selectivity means that the two overlapping transitions at one frequency (Sec.~\ref{sec:overlap})
saturate at different fractional rates. It can only arise from differences in $\omega_1$, the
matrix element, or $T_2$ between the two groups of spins, or from processes beyond single-spin
saturation.

\begin{enumerate}
\item \textbf{Orientation- and site-dependent $T_2$.} On resonance, $W\propto T_2$ at fixed
$\omega_1$. The homogeneous width comes from D--D and D--H dipolar couplings, from proximity
to paramagnetic centres, and from molecular motion. All of these depend on molecular
orientation and site (Sec.~\ref{sec:homogeneous}). This is the most direct physical route to
selectivity, and it can be measured from narrow-hole widths.

\item \textbf{The like-spin second moment.} In Van Vleck's theory \cite{VanVleck1948,Abragam1961},
a dipolar neighbour contributes $9/4$ times more to the second moment when it is ``like''
(resonant within the linewidth) than when it is ``unlike''. For quadrupole-split deuterons,
a neighbour counts as like only if its transition matches within the homogeneous width. The
homogeneous width at a frequency therefore grows with the local density of matching
transitions, by up to $(9/4)^{1/2}=1.5$ in width. This is a genuine density-dependent
per-spin effect. Its sign is to \emph{broaden} dense regions and so to lower their
on-resonance rate.

\item \textbf{Spatial $B_1$ inhomogeneity.} The coil, the target cup and the packing give a
distribution of $\omega_1$ across the sample. It is uncorrelated with frequency, so it
produces no selectivity. It does produce a slow, non-exponential tail of saturation (a
candidate contributor to apparent plateaus, Sec.~\ref{sec:open}).

\item \textbf{Hyperfine-shifted spins} near radicals are detuned and saturate differently
(Sec.~\ref{sec:barrier}).

\item \textbf{Collective effects.} Radiation damping and RF-induced changes of coil loading
depend on the magnetization at a frequency. For deuterons at the signal levels of polarized
targets they are expected to be small.
\end{enumerate}

\status{Eq.~\eqref{eq:rfrate} is established for weak drive. The orientation dependence of
$T_2$ in d-butanol and ND$_3$ is \textbf{not measured}. No first-principles mechanism makes the
rate per spin proportional to powder density.}

\subsection{Strong drive, rotating-frame relaxation and coherent effects}
When $\omega_1\gtrsim$ the local dipolar field, Eq.~\eqref{eq:golden} breaks down:
\begin{itemize}
\item \textbf{Off resonance:} the spin system reaches a rotating-frame spin temperature
\cite{Redfield1955,Goldman1970}.
\item \textbf{On resonance:} spin locking preserves the magnetization against saturation,
which then decays with $T_{1\rho}$.
\item \textbf{Pulsed operation:} nutation and quadrupolar echoes, which can invert populations
rather than equalize them (AFP is the adiabatic limit, Sec.~\ref{sec:afp}).
\end{itemize}
In polarized-target practice the RF level is set by the Q-meter circuit. Whether
$\omega_1T_2<1$ holds must be checked against a measured $B_1$ calibration.

\subsection{Coupling of off-resonance RF to the nuclear dipolar reservoir}
\label{sec:provotorov}
Off-resonance RF transfers energy between the Zeeman system and the nuclear dipolar (and
quadrupolar non-Zeeman) reservoir \cite{Provotorov1962,Goldman1970}.

\textbf{Single homogeneous line.} Take a line at offset $\Delta$ from the RF, with
rotating-frame Zeeman inverse temperature $\alpha$, dipolar inverse temperature $\beta$ and
local field $D$ (in frequency units). Provotorov's equations are
\begin{equation}
\dot\alpha=-W(\Delta)(\alpha-\beta)-\frac{\alpha-\alpha_L}{T_{1Z}},\qquad
\dot\beta=-W(\Delta)\frac{\Delta^2}{D^2}(\beta-\alpha)-\frac{\beta-\beta_L}{T_{1D}} .
\label{eq:provotorov}
\end{equation}

\textbf{Packet form.} For an inhomogeneous (Pake) line at the population level, take a carrier
at $\omega_c$ driving a transition at $\omega$. The current acquires a term that is first
order in the dipolar inverse temperature $\beta_d$:
\begin{equation}
J=W\Big[(n_{\rm lower}-n_{\rm upper})-\hbar\beta_d(\omega-\omega_c)\,\tfrac12(n_{\rm lower}+n_{\rm upper})\Big],
\qquad
C_d\,\dot\beta_d=\sum(\omega-\omega_c)J+\dots
\label{eq:rfdipolar}
\end{equation}
\begin{itemize}
\item Absorption below resonance takes the energy mismatch from the reservoir and cools it.
\item This is the nuclear analogue of dynamic cooling. It also makes RF hole shapes
asymmetric.
\end{itemize}

\status{The theory is established \cite{Provotorov1962,Goldman1970,AbragamGoldman1982}. The
heat capacity $C_d$ of the deuteron non-Zeeman reservoir in d-butanol, and its thermal contact
with the quadrupole-split Zeeman packets, are \textbf{unknown}. The first-order form requires
$\hbar\beta_d|\omega-\omega_c|\ll1$.}

%==========================================================================================
\section{Adiabatic fast passage}
\label{sec:afp}
A frequency or field sweep through a transition at rate $\dot\omega$ inverts it with
Landau--Zener efficiency
\begin{equation}
\epsilon=1-\exp\!\Big(-\frac{\pi\,\Omega_R^2}{2|\dot\omega|}\Big)=1-\exp\!\Big(-\frac{\pi\,\omega_1^2}{|\dot\omega|}\Big),
\qquad \Omega_R=2\,|\bra{a}\omega_1 I_x\ket{b}|=\sqrt2\,\omega_1 ,
\end{equation}
where $\Omega_R$ is the Rabi frequency of the isolated two-level transition (the other
transition of the same spins is far off resonance). This holds provided the sweep is fast
compared with $T_{1\rho}$ and slow compared with
$\omega_1$ \cite{Abragam1961,Slichter1990,GarwoodDelaBarre2001}.

\textbf{Effect at one frequency.} A sweep through frequency $R$ inverts both transitions
there. It swaps $n_{+1}\leftrightarrow n_0$ for the packet whose I$_+$ lies at $R$, and
$n_0\leftrightarrow n_{-1}$ for its mirror packet.

\textbf{Gain in $Q$.} Using $Q=1-3\avg{n_0}$, the change from inverting at $R$ is
\begin{equation}
\Delta Q_{\rm AFP}(R)=3\epsilon\Big[\mu_{\rm m}\big(p_0-p_{-1}\big)_{\rm m}-\mu_{\rm i}\big(p_{+1}-p_0\big)_{\rm i}\Big],
\label{eq:afpgain}
\end{equation}
where i is the packet whose I$_+$ lies at $R$ and m its mirror. The gain is positive where the
mirror packet is more than $(p_{+1}-p_0)/(p_0-p_{-1})\approx2.3$ times denser (at $P=0.5$).

Inversion moves twice the population that saturation does, so selective AFP windows can raise
$Q$ well beyond what selective saturation can reach. For d-butanol the beneficial windows are
the negative horn region and the outer positive region beyond the C--D edge.

\textbf{Limits of the idealization.} The adiabatic condition, relaxation during the sweep,
sweep edges, and a sweep passing through both transitions of the same spins (which gives a
$+1\leftrightarrow-1$ exchange) limit real efficiencies.

\status{Physics established. The selective-window efficiencies achievable in a polarized-target
Q-meter set-up are \textbf{not measured}.}

%==========================================================================================
\section{Spin-lattice relaxation}
\label{sec:t1}

\subsection{Quadrupolar relaxation of spin 1}
Fluctuations of the EFG (molecular motion and phonons) induce $\Delta m=\pm1$ transitions at
rate $W_1$ and $\Delta m=\pm2$ transitions at rate $W_2$ \cite{Abragam1961,Wasylishen2012}.
Linearizing about equilibrium,
\begin{equation}
\dot P=-(W_1+2W_2)(P-P_L),\qquad \dot Q=-3W_1(Q-Q_L),
\label{eq:t1pq}
\end{equation}
with the rates defined per transition in each direction. Vector and tensor polarization
therefore relax with \emph{different} times,
\[
T_{1Z}^{-1}=W_1+2W_2,\qquad T_{1Q}^{-1}=3W_1 ,
\]
and both depend on orientation through the spectral densities.

\subsection{Paramagnetic relaxation}
In doped targets, nuclear relaxation is dominated by the radicals. Nuclei near a centre relax
directly through the fluctuating hyperfine field \cite{Bloembergen1949,Blumberg1960}. The bulk
relaxes by spin diffusion towards those centres, in the diffusion-limited or rapid-diffusion
regimes of Khutsishvili and Blumberg, which gives $T_1^{-1}\propto c_e$
\cite{Blumberg1960,AbragamGoldman1978}. A treatment specific to polarized targets is given
in \cite{JiZhi1999}. Both mechanisms are $\Delta m=\pm1$ dominated.

\subsection{Lattice and spin-temperature equilibria}
\label{sec:equilibrium}
At the lattice temperature $T_L\simeq1$\,K, $P_L\approx 10^{-3}$. Under DNP, the nuclear
Zeeman system instead approaches a spin temperature set by the electrons
(Sec.~\ref{sec:dnp}).

\textbf{Quadrupole corrections.} The full Boltzmann distribution includes the quadrupolar
energy, $n_m\propto\exp[-\beta E_m]$. This gives orientation-dependent corrections to
Eq.~\eqref{eq:boltzmann} of relative order $\beta\hbar\omega_Q$. At deuteron spin temperatures
corresponding to $P\simeq0.5$ these are about $0.4\%$.

\textbf{Separate quadrupolar order.} The tensor (quadrupolar-order) part can acquire its own
temperature after a manipulation. It returns to the value set by $P$ only through processes
that change $n_0$: lattice relaxation ($3W_1$), cross-branch spin flips
(Sec.~\ref{sec:diffusion}), and DNP.

\status{Mechanisms established. Orientation dependence, the $W_2/W_1$ ratio and the absolute
$T_1$ of deuterons in d-butanol at 1\,K are \textbf{not separately measured}; in practice
$T_1$ is hours to days.}

%==========================================================================================
\section{Spectral spin diffusion for spin 1}
\label{sec:diffusion}

\subsection{Dipolar flip-flops between quadrupole-split spins}
Two deuterons $i,j$ couple through the secular dipolar Hamiltonian
\cite{Abragam1961,Bloembergen1949}
\[
\mathcal H_{ij}=b_{ij}\big[2I_{zi}I_{zj}-\tfrac12(I_{+i}I_{-j}+I_{-i}I_{+j})\big],\qquad
b_{ij}=\frac{\mu_0\gamma_d^2\hbar}{4\pi r_{ij}^3}\,\tfrac12(1-3\cos^2\vartheta_{ij}),
\]
where $\vartheta_{ij}$ is the angle of the internuclear vector to $\mathbf B_0$. For spin 1,
\[
\bra{+1}I_+\ket{0}=\bra{0}I_+\ket{-1}=\sqrt2 .
\]
The flip-flop term therefore connects the following pair states, each with a matrix element of
magnitude $|b_{ij}|$:

\begin{center}
\begin{tabular}{@{}lll@{}}
\toprule
Channel & Pair transition & Energy matched when \\
\midrule
same-branch I$_+$ & $\ket{0_i,{+1}_j}\leftrightarrow\ket{{+1}_i,0_j}$ & $\omega_+(\Omega_i)=\omega_+(\Omega_j)$\\
same-branch I$_-$ & $\ket{{-1}_i,0_j}\leftrightarrow\ket{0_i,{-1}_j}$ & $\omega_-(\Omega_i)=\omega_-(\Omega_j)$\\
cross-branch & $\ket{0_i,0_j}\leftrightarrow\ket{{+1}_i,{-1}_j}$ & $\omega_+(\Omega_i)=\omega_-(\Omega_j)$\\
\bottomrule
\end{tabular}
\end{center}

\textbf{Same matrix element, different partners.} All three channels have the same dipolar
matrix element. The cross-branch process exchanges excitation between the two
\emph{overlapping} lines at one spectral position: it couples the I$_+$ of one spin with the
I$_-$ of a spin at the mirror orientation. In a random powder its strength relative to the
same-branch channels is set by
\begin{itemize}
\item the density of energy-matched partners: same-branch partners at frequency $R$ have
capacity $\mu(R)$, cross-branch partners the mirror capacity;
\item any orientation correlation between dipolar-coupled neighbours (intramolecular
deuterons have fixed relative bond orientations).
\end{itemize}
A cross-branch/same-branch coupling ratio of order unity is therefore the \emph{expected}
value for uncorrelated pairs, not an arbitrary choice.

\textbf{Double-quantum channel.} The exchange
$\ket{{+1}_i,{-1}_j}\leftrightarrow\ket{{-1}_i,{+1}_j}$ needs two flip-flops. To first order
it conserves both Zeeman and quadrupolar energy for \emph{any} orientations, because $E_{+1}$
and $E_{-1}$ have the same quadrupole shift. It is therefore nonlocal in the spectrum. Through
an off-resonant intermediate $\ket{0,0}$, its amplitude is of order $b_{ij}^2/\delta$, where
$\delta$ is the intermediate energy mismatch, so its rate is suppressed by roughly
$(b/\delta)^2$ relative to resonant flip-flops.

\textbf{Quadrupolar quenching.} Flip-flops between spins whose transitions differ by more than
the homogeneous width are truncated. Spectral diffusion in resolved quadrupolar spectra is
therefore local in frequency and channel-specific \cite{SuterErnst1985}. Energy mismatches can
be bridged by the nuclear dipolar reservoir (three-spin processes, Sec.~\ref{sec:thermalmix}).

\subsection{Pair-current rate equations}
Let $S(\Delta)$ be the zero-quantum spectral overlap of the two transitions involved (width of
order the homogeneous width or the dipolar local field), and let $C_{ij}$ describe spatial and
orientation correlation. The packet equations are reversible pair currents.

\textbf{Same-branch (I$_+$ shown; I$_-$ analogous):}
\begin{equation}
J^{+}_{ij}=K\,\mu_i\mu_j\,C_{ij}\,S(R_i-R_j)\big(p_{0i}p_{+j}-p_{+i}p_{0j}\big).
\end{equation}

\textbf{Cross-branch} (the I$_+$ of $i$ at $R_i$ with the I$_-$ of $j$ at $-R_j$):
\begin{equation}
J^{X}_{ij}=K r_X\,\mu_i\mu_j\,C_{ij}\,S(R_i+R_j)\big(p_{0i}p_{0j}-p_{+i}p_{-j}\big).
\end{equation}

\textbf{Double-quantum:}
\begin{equation}
J^{DQ}_{ij}=K r_{DQ}\,\mu_i\mu_j\,C_{ij}\big(p_{-i}p_{+j}-p_{+i}p_{-j}\big).
\end{equation}

\textbf{Properties} \cite{SuterErnst1985,AbragamGoldman1982}:
\begin{itemize}
\item All currents conserve packet mass and total $P$.
\item Same-branch currents conserve each level total separately.
\item Cross-branch currents change $Q$ at fixed $P$; they are a spin-spin route by which
tensor order relaxes to the value set by $P$.
\item Double-quantum currents move vector polarization across the spectrum without changing
$Q$.
\item Every current vanishes for a common Boltzmann distribution. The stationary state is
unique when the exchange graph is connected.
\end{itemize}
The currents are bilinear in the populations. Spin diffusion therefore slows where partner
transitions are saturated ("flip-flop blocking"), and it cannot be represented as each packet
relaxing to a fixed reference.

\subsection{Continuum limit: diffusion along the Pake line}
\label{sec:continuum}
For small spectral jumps, the kernel reduces to a diffusion equation in frequency. Linearizing
for a single branch with $\mu(\omega)$ the Pake density,
\begin{equation}
\partial_t p=\frac{1}{\mu(\omega)}\,\partial_\omega\!\big[\mu(\omega)\,\tilde D(\omega)\,\partial_\omega p\big],
\qquad \mu(\omega)\propto\frac{\sin\theta}{|d\omega/d\theta|}.
\label{eq:continuum}
\end{equation}
This is the earlier draft's Laplace--Beltrami operator
$\sin^{-1}\theta\,\partial_\theta(D\sin\theta\,\partial_\theta n)$ rewritten with the
Jacobian $J=|d\omega/d\theta|=\Omega\sqrt{2(1-x)(1+2x)}$, $x=\omega/\Omega$, and
$\tilde D=DJ^2$.

Three caveats apply:
\begin{enumerate}
\item The variable $\theta$ labels packets. Spins do not reorient; ``diffusion in $\theta$''
describes energy-conserving exchange between molecules.
\item Cross-branch and double-quantum transport are nonlocal and \emph{cannot} be written in
this form.
\item The no-flux boundary conditions apply at the edges of the support,
$\omega\in\{-\Omega/2,\Omega\}$.
\end{enumerate}

\subsection{Real space, diffusion barrier and microwaves}
\label{sec:barrier}
Spectral transport requires spatial proximity, $|b_{ij}|\propto r_{ij}^{-3}$. A Pake histogram
does not specify which spectral packets are spatial neighbours: the correlation $C_{ij}$ is an
effective closure.

\textbf{Diffusion barrier.} Deuterons within the "diffusion barrier" of a radical, where the
hyperfine shift exceeds the dipolar coupling, are detuned from the bulk and exchange with it
only weakly \cite{Khutsishvili1966,AbragamGoldman1978,Wenckebach2016}.

\textbf{Microwave effects.} Microwave irradiation can modify nuclear spin diffusion by reducing
the electron polarization. Electron-polarization-dependent proton spin diffusion has been
observed at 1.2\,K \cite{Chessari2023}.

\status{The mechanisms are established. The deuteron K, $S$ width, $r_X$, $r_{DQ}$ and
correlations in d-butanol are \textbf{weakly identified}. A fit to one burn/refill experiment
gave a nonlocal channel much stronger than the $(b/\delta)^2$ estimate. That channel is
therefore effective, and may absorb reservoir-assisted or real-space transport.}

%==========================================================================================
\section{Thermal mixing among the nuclei}
\label{sec:thermalmix}
Flip-flops between packets of different frequency are allowed when the energy mismatch is
supplied or absorbed by the nuclear dipolar reservoir (three-spin processes)
\cite{Goldman1970,AbragamGoldman1982}. With reservoir inverse temperature $\beta_d$, detailed
balance gives, to first order,
\begin{equation}
J_{ij}=H_{ij}\Big[(p_{0i}p_{+j}-p_{+i}p_{0j})+\tfrac12\hbar\beta_d\,E_{ij}\,(p_{0i}p_{+j}+p_{+i}p_{0j})\Big],
\label{eq:mixing}
\end{equation}
where $E_{ij}$ is the energy released. Analogous terms apply to the cross-branch channel, with
$E=\hbar(R_i+R_j)$. The double-quantum channel releases no energy at first order.

\textbf{Energy conservation.} The spin offset energy
$E_{\rm spin}=-\hbar\sum R(n_{+1}+n_{-1})$ and the reservoir energy $E_d=-C_d\beta_d$ are
conserved together:
\[
\dot E_{\rm spin}+\dot E_d=0 .
\]
Thermal mixing drives the whole line towards a common rotating-frame spin temperature. It
broadens the effective exchange width beyond the homogeneous width, and it couples to RF
through Eq.~\eqref{eq:rfdipolar}.

\status{Theory established. Its importance for deuterons in d-butanol, and $C_d$, are
\textbf{unknown}.}

%==========================================================================================
\section{Dynamic nuclear polarization of spin-1 nuclei}
\label{sec:dnp}

\subsection{Overview and regimes}
Microwaves at $\omega_\mu$ near the electron Larmor frequency $\omega_e$ (140\,GHz at 5\,T)
transfer electron polarization to the nuclei. Which mechanism dominates depends on the ESR
width $\Delta_e$ relative to $\omega_d$, on electron spectral diffusion and on $c_e$
\cite{Jeffries1963,AbragamGoldman1978,deBoer1976,Wenckebach2016}:

\begin{center}
\begin{tabular}{@{}>{\raggedright\arraybackslash}p{3.4cm}>{\raggedright\arraybackslash}p{5.2cm}>{\raggedright\arraybackslash}p{6.6cm}@{}}
\toprule
Regime & Condition & Mechanism \\
\midrule
Solid effect (SE) & $\Delta_e\ll\omega_d$ & forbidden e--n flips at $\omega_e\pm\omega_n$\\
Differential SE (DSE) & $\Delta_e\gtrsim\omega_d$, weak e--e spectral diffusion & overlapping SE sidebands, net difference\\
Cross effect (CE) & $\Delta_e>\omega_d$, e--e coupled pairs & three-spin $e_1e_2n$ with $|\omega_{e1}-\omega_{e2}|=\omega_n$\\
Thermal mixing (TM) & $\Delta_e>\omega_d$, fast e--e spectral diffusion & spin-temperature contact with electron dipolar reservoir\\
\bottomrule
\end{tabular}
\end{center}

For d-butanol with TEMPO or EHBA-Cr(V) at 5\,T, $\Delta_e$ is several hundred MHz, far larger
than $\omega_d/2\pi=32.7$\,MHz. DNP is then usually described by spin-temperature theory (TM),
with DSE contributions discussed for ammonia
\cite{Goertz2004,Meyer2004,deBoer1976,Wenckebach2016}.

For spin 1 every mechanism acts \emph{separately} on the two transitions of each orientation,
at their own frequencies $\omega_\pm(\Omega)$.

\subsection{Solid effect}
The hyperfine interaction mixes states, so microwaves at
\[
\omega_\mu\simeq\omega_e\mp\omega_\pm(\Omega)
\]
flip the electron and the nucleus together \cite{Jeffries1963,AbragamGoldman1978,Hovav2010}.
In the weak-mixing limit the effective operator is
\[
\mathcal H^{\rm eff}_{SE}\propto(A_\perp/\omega_d)\,(S_+I_-+S_-I_+),
\]
so the transition rate is proportional to $(A_\perp/\omega_d)^2\omega_{1e}^2$ times the
electron spectral density. (The mixing coefficient involves $\omega_n$ in the denominator, not
$\omega_e$.)

For each branch $b\in\{+,-\}$, nucleus-raising ($\uparrow$) and lowering ($\downarrow$) rates
are
\begin{equation}
W_b^{\uparrow,\downarrow}(\Omega;\omega_\mu)=k_{SE}\,c_e\avg{A_\perp^2}\,
\frac{\omega_{1e}^2T_{2e}}{1+\Delta_{\mp}^2T_{2e}^2+\omega_{1e}^2T_{1e}T_{2e}}\,
L_e\big(\Delta_\mp\big),\qquad
\Delta_\mp=\omega_\mu-\big[\omega_e\mp\omega_b(\Omega)\big].
\label{eq:se}
\end{equation}
Here $L_e$ is the ESR line and $\beta=\omega_{1e}^2T_{1e}T_{2e}/(1+\Delta^2T_{2e}^2)$ is
Jeffries' pumping parameter.

The net pumping direction is proportional to the local electron polarization,
$P_e^{\rm loc}=P_e^{(0)}/(1+\beta)$. Double-quantum SE ($\Delta m_I=\pm2$) is allowed only at
higher order and is weak.

The spin-1 SE rate equations in the narrow-line limit, including phonon channels, are given in
\cite{FeddersSouers1987}. The orientation-resolved generalization adds three things:
\begin{itemize}
\item the full $\Omega$ dependence through $\omega_\pm(\Omega)$;
\item arbitrary ESR shape;
\item spin diffusion (Sec.~\ref{sec:diffusion}).
\end{itemize}

\subsection{Differential solid effect}
When $\Delta_e\gtrsim\omega_d$, the positive and negative SE sidebands from different spin
packets overlap. The net polarization comes from the \emph{difference} between packets on the
two sides of $\omega_\mu$ \cite{Wenckebach2016}. The rates become integrals of
Eq.~\eqref{eq:se} over the ESR line:
\begin{equation}
W_b^{\uparrow}-W_b^{\downarrow}\propto
\int d\omega\,s_e(\omega-\omega_\mu)\,L_e(\omega)\,
\big[L_e(\omega+\omega_b)-L_e(\omega-\omega_b)\big].
\end{equation}
The result is a smooth S-curve whose zero crossing is shifted by line asymmetry
($g$-anisotropy) and by saturation gradients.

\subsection{Cross effect}
Two dipolar-coupled electrons whose Larmor frequencies differ by a nuclear transition
frequency,
\[
|\omega_{e1}-\omega_{e2}|\simeq\omega_\pm(\Omega),
\]
exchange polarization with the nucleus through the three-spin term
$S_{1\pm}S_{2\mp}I_{\pm}$ \cite{Hovav2012,Wenckebach2016}. Microwaves saturate one electron of
the pair. The rate scales with the pair density, and so with $c_e^2$:
\begin{equation}
W^{CE,\uparrow\downarrow}_b\propto c_e^2\avg{b_{ee}^2}\!\int\! d\omega\,s_e(\omega-\omega_\mu)L_e(\omega)L_e(\omega\mp\omega_b)
\big[P_e(\omega\mp\omega_b)-P_e(\omega)\big].
\end{equation}

\subsection{Thermal mixing and spin temperature}
\label{sec:tm}
If electron spectral diffusion is fast, the electron non-Zeeman (dipolar plus offset) energy
has a single inverse temperature $\beta$ in the frame rotating at $\omega_\mu$
\cite{Provotorov1962,Borghini1968,AbragamGoldman1978,Goldman1970}.

\textbf{Homogeneous ESR line.} With offset $\Delta=\omega_e-\omega_\mu$, Provotorov's equations
\eqref{eq:provotorov} apply to the electrons. Under strong saturation ($\alpha=\beta$) the
steady state is
\begin{equation}
\beta\simeq\beta_L\,\frac{\omega_e\,\Delta}{\Delta^2+D_e^2\,T_{1Z}/T_{1D}} .
\label{eq:tmhomog}
\end{equation}
This is odd in $\Delta$, with extrema at $|\Delta|=D_e\sqrt{T_{1Z}/T_{1D}}$, where $D_e$ is the
electron local field.

\textbf{Inhomogeneous ESR line.} Borghini's spin-temperature model gives the electron
polarization profile across the line, and $\beta$ is fixed by energy balance between
microwave absorption and electron spin-lattice relaxation \cite{Borghini1968,AbragamGoldman1978}.

\textbf{Contact with the nuclei.} Through three-spin e--e--n flips, the nuclear Zeeman system
approaches the same $\beta$:
\begin{equation}
\dot{\bm n}(\Omega)\big|_{TM}=-\frac{1}{T_{IS}(\Omega)}\Big[\bm n(\Omega)-\bm n^{ST}(\beta;\Omega)\Big],\qquad
n^{ST}_{+1}:n^{ST}_0:n^{ST}_{-1}=e^{\beta\hbar\omega_+}:1:e^{-\beta\hbar\omega_-},
\label{eq:tmcontact}
\end{equation}
with each transition driven toward the ratio set by its own frequency. The thermal-contact time
satisfies $T_{IS}^{-1}\propto c_e$ times the electron dipolar spectral density at the nuclear
frequency.

Two consequences:
\begin{enumerate}
\item DNP by TM drives $P$ and $Q$ toward the spin-temperature relation \eqref{eq:boltzmann},
up to quadrupole corrections of order $\beta\hbar\omega_Q$. It acts as an \emph{absolute}
target, $\beta(\omega_\mu)$, not as relaxation toward the current $P$.
\item If the contact is stronger near the horns or for particular sites, DNP itself is
orientation selective.
\end{enumerate}

\textbf{Heat balance.} A consistent three-reservoir model (electron Zeeman, electron dipolar,
nuclear Zeeman, inverse temperatures $\alpha$, $\beta$, $\beta_n$) has the form
\begin{equation}
C_n\dot\beta_n=-\frac{C_n(\beta_n-\beta)}{T_{IS}}-\frac{C_n(\beta_n-\beta_L)}{T_{1n}},\qquad
C_D\dot\beta=\dots+\frac{C_n(\beta_n-\beta)}{T_{IS}},
\end{equation}
where the back-reaction of the nuclei on the electron dipolar reservoir ("nuclear heat load")
matters for high nuclear heat capacity.

\textbf{Microwave on/off asymmetry.} Microwaves can also change nuclear transport (by reducing
the electron polarization) separately from pumping \cite{Chessari2023}. An observation that
about three times more RF power is needed to reach the same peak residual with microwaves on
constrains the sum of these effects. It does not separate them.

\status{SE, CE, TM and DSE theory are established in the cited texts. For spin-1 deuterons in
d-butanol the relative weight of the mechanisms, and their orientation dependence, are
\textbf{not identified}. A phenomenological ``relax to the spin-temperature target'' form, with
rate and saturation $P$ fitted to build-up data, describes the measured 2014 build-up to 0.8\%
RMS.}

%==========================================================================================
\section{The complete master equation}
\label{sec:master}
Collect all processes for packet $k=(R,s)$. Its I$_+$ transition lies at $R$ and its I$_-$ at
the mirror.
\begin{equation}
\boxed{\;
\dot{\bm n}_k=\underbrace{\bm R_{\rm RF}(\bm U,\beta_d)\bm n_k}_{\rm Sec.~\ref{sec:rf}}
+\underbrace{\bm R_{\rm DNP}\bm n_k+\bm s^{\rm DNP}_k}_{\rm Sec.~\ref{sec:dnp}}
+\underbrace{\bm R_{T_1}(\bm n_k-\bm n^{L}_k)}_{\rm Sec.~\ref{sec:t1}}
+\underbrace{\sum_{l}\bm J_{kl}(\bm n,\beta_d)}_{\rm Secs.~\ref{sec:diffusion},\ref{sec:thermalmix}}
\;}
\label{eq:master}
\end{equation}
Here:
\begin{itemize}
\item $\bm R_{\rm RF}$ has rates
$g^{\pm}_k=c_k\sum_j\kappa\Gamma_k/[\Gamma_k^2+(\omega^\pm_k-\omega_j)^2]\,U_j$, where $U_j$ is
the carrier power, $c_k$ the per-spin coupling and $\Gamma_k=1/T_{2,k}$.
\item $\bm R_{T_1}$ contains $W_1$ and $W_2$.
\item $\bm J_{kl}$ are the three spin-diffusion channels, with the thermal-mixing corrections of
Eq.~\eqref{eq:mixing}.
\item With the dipolar reservoir included, $\beta_d$ obeys
$C_d\dot\beta_d=\sum({\rm offset})\times({\rm current})-C_d(\beta_d-\beta_L)/T_{1d}$.
\end{itemize}
AFP enters as an instantaneous map between integration intervals.

\paragraph{Conservation and balance.}
\begin{itemize}
\item Every term preserves packet mass.
\item Spin diffusion conserves $P$.
\item Exchange together with the reservoir conserves offset energy.
\item The RF terms remove $P$ and drive towards local saturation.
\item DNP and $T_1$ are the only sources that restore $P$.
\end{itemize}
Without RF, the unique stationary state is the spin-temperature state of the dominant source,
modified by the reservoir.

\paragraph{Discretization.} The spectral grid must resolve the narrowest of $\Gamma_k$, the
exchange width and the line edges. Pair sums over Lorentzian tails are Toeplitz (or Hankel, for
the cross-branch channel) and are evaluated efficiently by FFT. Forward-Euler integration must
respect $h\sum(\text{outgoing rates})<1$.

%==========================================================================================
\section{General consequences for tensor polarization}
\label{sec:consequences}
The following results follow from Eqs.~\eqref{eq:PQ}, \eqref{eq:rfrate} and \eqref{eq:master}.
The numbers in brackets are for d-butanol at $P_0=0.50$ ($Q_0=0.197$); they come from numerical
optimization in the fitted model, are model-dependent, and are given only as orders of
magnitude.

\begin{enumerate}
\item \textbf{Saturation versus overlap.} Saturating the I$_-$ transition of a packet with
capacity $\mu$ raises $Q$ by $\tfrac32\mu(p_0-p_{-1})$ [$0.151\mu$]. Saturating an I$_+$
lowers it by $\tfrac32\mu(p_{+1}-p_0)$ [$0.349\mu$]. A saturating burn at one frequency always
affects both overlapping transitions, so its net value depends on their capacities and on
their relative per-spin rates.

\item \textbf{Equal per-spin rates.} The only lever left is the choice of frequencies.
Optimized single burns reach [$Q\approx0.237$].

\item \textbf{Perfect selectivity ceiling.} The two transitions of one packet share its
coupling, so at most the denser packet of each mirror pair can be selectively saturated:
\[
Q\le Q_0+\tfrac32(p_0-p_{-1})\sum_{\rm pairs}\max(\mu_k,\mu_{\bar k})
\]
[0.360]. If the sparser packet is favoured instead, the bound is [0.261]. The bound for
saturating all I$_-$ transitions with no overlap at all, $Q_0+\tfrac32(p_0-p_{-1})$ [0.424],
is unreachable.

\item \textbf{Inversion.} AFP gains are twice the saturation gains at the same spectral
position, Eq.~\eqref{eq:afpgain}. Two-window AFP gives [0.298 at $\epsilon=1$; 0.278 at
$\epsilon=0.8$], independent of the unknown RF selectivity.

\item \textbf{Per-spin selectivity sets most of the single-burn range}
[$0.232$--$0.272$ over plausible to extreme selectivity]. Refilling strength and time-shaping
add less [$\lesssim0.009$]. Site resolution and the dipolar reservoir change single-burn optima
by less than $10^{-3}$.

\item \textbf{Steady state with continuous DNP.} DNP restores $P$ while a weak burn converts
it to tensor order. The optimal constant burn sits where only the I$_-$ line exists (beyond the
C--D edge). It gives [$Q=0.202$--$0.213$, depending on the transport closure]. Burning the horn
at the measured DNP rate lowers $Q$ [0.175], because DNP is too slow (435\,s) to replace $P$.

\item \textbf{Transport type is decisive.} If spectral transport stayed within each branch,
the I$_-$ saturation produced in the I$_-$-only region would be carried into the overlap
region while DNP resupplies $P$ [steady $Q$ up to 0.24--0.30]. With cross-branch transport
of comparable strength, as the equal matrix elements of Sec.~\ref{sec:diffusion} suggest for
uncorrelated pairs, stronger transport \emph{lowers} $Q$ [0.18 and 0.15 at 10$\times$ and
100$\times$]. This mechanism is the largest open lever, and it is constrained by measurable
refill behaviour.

\item \textbf{Timing.} For incoherent saturation the response to RF dose is concave:
already-equalized transitions gain nothing from more power. Spreading a given dose evenly in
time is optimal. Pulse trains, sweeps and optimized periodic schedules did not beat a steady
burn [gain $\le0.0016$; long gaps lose up to 0.024]. A transient overshoot after RF switch-on
[0.003--0.009] can be used once. \emph{This conclusion does not extend to coherent high-power
sequences}, which can invert populations.
\end{enumerate}

%==========================================================================================
\section{What is known and what is not}
\label{sec:open}

\begin{longtable}{@{}p{3.6cm}p{3.6cm}p{3.1cm}p{4.6cm}@{}}
\toprule
Mechanism & Theory & d-butanol status & Measurement that would pin it\\
\midrule
\endhead
Pake spectrum, sites & established \cite{Wasylishen2012,Pake1948} & fitted (0.8\%) & --\\
Weak-drive RF rate & established \cite{Abragam1961} & strength fitted to one burn & $B_1$ calibration; power/duration series\\
Per-spin selectivity ($T_2(\Omega)$) & mechanisms known \cite{VanVleck1948} & \textbf{unknown} & narrow-hole widths and initial rates at horn vs pedestal; mirror response\\
$B_1$ inhomogeneity & known & \textbf{unknown} & nutation; saturation-curve shape\\
AFP efficiency & established \cite{GarwoodDelaBarre2001} & \textbf{unknown} for selective windows & selective sweeps with line-shape analysis\\
$T_1$: $W_1$, $W_2$ & established \cite{Abragam1961} & \textbf{not separated} & relaxation of $P$ and $Q$ after manipulation\\
Same-branch diffusion ($K$, width) & established \cite{SuterErnst1985} & weakly identified & refill after narrow burns\\
Cross-branch ratio $r_X$ & expected $\sim1$ (equal matrix elements) & \textbf{never fitted} & refill of I$_+$ vs I$_-$ after a burn in the I$_-$-only region\\
DQ exchange & second order, weak & effective channel fitted strong & refill far from the burn\\
Nuclear dipolar reservoir & established \cite{Provotorov1962,Goldman1970} & \textbf{unknown} $C_d$ & hole asymmetry; off-resonance RF\\
DNP mechanism and selectivity & established \cite{AbragamGoldman1978,Wenckebach2016} & effective rate fitted & frequency-resolved build-up after burns\\
Microwave-modified transport & observed for protons \cite{Chessari2023} & \textbf{unknown} & RF burn/refill with microwaves on vs off\\
\bottomrule
\end{longtable}

\paragraph{Open experimental puzzles.}
\begin{enumerate}
\item \textbf{The persistent $\sim20\%$ RF-on peak residual.} At about 0.1\,W the peak falls
quickly, then stalls near 20\% under continued irradiation, while the pedestal burns to about
zero; higher power burns the peak further. No fitted rate model reproduces this. Candidate
explanations:
\begin{itemize}
\item spatial $B_1$ inhomogeneity (a weakly driven fraction);
\item strong replenishment of the horn from dense same-branch partners;
\item hyperfine-shifted or otherwise detuned sub-ensembles;
\item rotating-frame or dipolar-reservoir saturation, which is not pure equalization;
\item RF heating.
\end{itemize}
Plateau duration and power dependence would discriminate between them.

\item \textbf{Microwave on/off power ratio} (about $3\times$). This requires separating DNP
pumping from microwave-modified transport.

\item \textbf{Deviations from equal spin temperature} between $P$ and $Q$ under DNP in
ND$_3$ \cite{Keller2017}. This tests whether quadrupolar order has its own temperature or its
own contact to the electrons.
\end{enumerate}

%==========================================================================================
\appendix
\section{Spin-1 operators}
In the basis $(\ket{+1},\ket{0},\ket{-1})$:
\[
I_z=\begin{pmatrix}1&0&0\\0&0&0\\0&0&-1\end{pmatrix},\quad
I_+=\sqrt2\begin{pmatrix}0&1&0\\0&0&1\\0&0&0\end{pmatrix},\quad
I_x=\tfrac12(I_++I_-),\quad 3I_z^2-2={\rm diag}(1,-2,1).
\]
The tensor polarization is $Q=\avg{3I_z^2-2}$ and the vector polarization is $P=\avg{I_z}$.

\section{Pake Jacobian in frequency coordinates}
With $\omega=\Omega\Pt(\cos\theta)$, we have $\cos^2\theta=(1+2\omega/\Omega)/3$ and
$\sin^2\theta=\tfrac23(1-\omega/\Omega)$. The Jacobian is
\[
J=|d\omega/d\theta|=\Omega\sqrt{2(1-\omega/\Omega)(1+2\omega/\Omega)},
\]
and the density is $\mu\propto\sin\theta/J\propto(1+2\omega/\Omega)^{-1/2}$, consistent with
Eq.~\eqref{eq:pakedensity}. Equation~\eqref{eq:continuum} then follows from the
Laplace--Beltrami form, with no-flux boundaries at $\omega=-\Omega/2$ and $\omega=\Omega$.

\section{Rate-equation matrices}
For one packet, write $u_\pm$ for the rate that raises $m$ across the I$_\pm$ transition and
$d_\pm$ for the rate that lowers it:
\[
\bm R=\begin{pmatrix}-d_+&u_+&0\\ d_+&-u_+-d_-&u_-\\0&d_-&-u_-\end{pmatrix},
\]
For RF, $u_\pm=d_\pm=g_\pm$; for DNP, $u\neq d$. The relaxation part to the lattice is
\[
\bm R_{T_1}=\begin{pmatrix}-W_1-W_2&W_1&W_2\\W_1&-2W_1&W_1\\W_2&W_1&-W_1-W_2\end{pmatrix}
\]
(high-temperature form; detailed-balance factors $e^{\pm\beta_L\hbar\omega}$ multiply the
upward rates).

%==========================================================================================
\begin{thebibliography}{99}

\bibitem{Abragam1961}
A.~Abragam, \emph{The Principles of Nuclear Magnetism} (Clarendon Press, Oxford, 1961).

\bibitem{Slichter1990}
C.~P.~Slichter, \emph{Principles of Magnetic Resonance}, 3rd ed. (Springer, Berlin, 1990).

\bibitem{Wasylishen2012}
R.~E.~Wasylishen, S.~E.~Ashbrook, and S.~Wimperis (eds.), \emph{NMR of Quadrupolar Nuclei in
Solid Materials} (John Wiley \& Sons, Chichester, 2012).

\bibitem{CohenReif1957}
M.~H.~Cohen and F.~Reif, ``Quadrupole effects in nuclear magnetic resonance studies of
solids,'' Solid State Physics \textbf{5}, 321 (1957).

\bibitem{Pake1948}
G.~E.~Pake, ``Nuclear resonance absorption in hydrated crystals: fine structure of the proton
line,'' J.\ Chem.\ Phys.\ \textbf{16}, 327 (1948).

\bibitem{VanVleck1948}
J.~H.~Van~Vleck, ``The dipolar broadening of magnetic resonance lines in crystals,''
Phys.\ Rev.\ \textbf{74}, 1168 (1948).

\bibitem{Bloembergen1949}
N.~Bloembergen, ``On the interaction of nuclear spins in a crystalline lattice,''
Physica \textbf{15}, 386 (1949).

\bibitem{Redfield1955}
A.~G.~Redfield, ``Nuclear magnetic resonance saturation and rotary saturation in solids,''
Phys.\ Rev.\ \textbf{98}, 1787 (1955).

\bibitem{Provotorov1962}
B.~N.~Provotorov, ``Magnetic resonance saturation in crystals,'' Sov.\ Phys.\ JETP
\textbf{14}, 1126 (1962).

\bibitem{Goldman1970}
M.~Goldman, \emph{Spin Temperature and Nuclear Magnetic Resonance in Solids} (Clarendon Press,
Oxford, 1970).

\bibitem{AbragamGoldman1978}
A.~Abragam and M.~Goldman, ``Principles of dynamic nuclear polarisation,'' Rep.\ Prog.\ Phys.\
\textbf{41}, 395 (1978).

\bibitem{AbragamGoldman1982}
A.~Abragam and M.~Goldman, \emph{Nuclear Magnetism: Order and Disorder} (Clarendon Press,
Oxford, 1982).

\bibitem{Jeffries1963}
C.~D.~Jeffries, \emph{Dynamic Nuclear Orientation} (Interscience, New York, 1963).

\bibitem{Borghini1968}
M.~Borghini, ``Spin-temperature model of nuclear dynamic polarization using free radicals,''
Phys.\ Rev.\ Lett.\ \textbf{20}, 419 (1968).

\bibitem{Blumberg1960}
W.~E.~Blumberg, ``Nuclear spin-lattice relaxation caused by paramagnetic impurities,''
Phys.\ Rev.\ \textbf{119}, 79 (1960).

\bibitem{Khutsishvili1966}
G.~R.~Khutsishvili, ``Spin diffusion,'' Sov.\ Phys.\ Usp.\ \textbf{8}, 743 (1966).

\bibitem{SuterErnst1985}
D.~Suter and R.~R.~Ernst, ``Spin diffusion in resolved solid-state NMR spectra,''
Phys.\ Rev.\ B \textbf{32}, 5608 (1985).

\bibitem{deBoer1976}
W.~de~Boer, ``Dynamic orientation of nuclei at low temperatures,'' J.\ Low Temp.\ Phys.\
\textbf{22}, 185 (1976).

\bibitem{Hovav2010}
Y.~Hovav, A.~Feintuch, and S.~Vega, ``Theoretical aspects of dynamic nuclear polarization in
the solid state -- the solid effect,'' J.\ Magn.\ Reson.\ \textbf{207}, 176 (2010).

\bibitem{Hovav2012}
Y.~Hovav, A.~Feintuch, and S.~Vega, ``Theoretical aspects of dynamic nuclear polarization in
the solid state -- the cross effect,'' J.\ Magn.\ Reson.\ \textbf{214}, 29 (2012).

\bibitem{Wenckebach2016}
W.~T.~Wenckebach, \emph{Essentials of Dynamic Nuclear Polarization} (Spindrift Publications,
The Netherlands, 2016).

\bibitem{GarwoodDelaBarre2001}
M.~Garwood and L.~DelaBarre, ``The return of the frequency sweep: designing adiabatic pulses
for contemporary NMR,'' J.\ Magn.\ Reson.\ \textbf{153}, 155 (2001).

\bibitem{CrabbMeyer1997}
D.~G.~Crabb and W.~Meyer, ``Solid polarized targets for nuclear and particle physics
experiments,'' Annu.\ Rev.\ Nucl.\ Part.\ Sci.\ \textbf{47}, 67 (1997).

\bibitem{Goertz2004}
St.~Goertz, ``The dynamic nuclear polarization process,'' Nucl.\ Instrum.\ Methods A
\textbf{526}, 28 (2004).

\bibitem{Meyer2004}
W.~Meyer, ``Ammonia as a polarized solid target material -- a review,'' Nucl.\ Instrum.\
Methods A \textbf{526}, 12 (2004).

\bibitem{Dulya1997}
C.~Dulya \emph{et al.}, ``A line-shape analysis for spin-1 NMR signals,'' Nucl.\ Instrum.\
Methods A \textbf{398}, 109 (1997).

\bibitem{FeddersSouers1987}
P.~A.~Fedders and P.~C.~Souers, ``Solid-effect rate equations for a spin-1 nucleus,''
Phys.\ Rev.\ B \textbf{35}, 3088 (1987).

\bibitem{JiZhi1999}
L.~Ji~Zhi, ``Spin relaxation in polarized targets with paramagnetic centers,''
Commun.\ Theor.\ Phys.\ \textbf{31}, 619 (1999).

\bibitem{Chessari2023}
A.~Chessari \emph{et al.}, Phys.\ Rev.\ B \textbf{107}, 224429 (2023).

\bibitem{KellerSSRF}
D.~Keller, ``Selective semi-saturation measurement of spin-1,'' Eur.\ Phys.\ J.\ A \textbf{62}, 36 (2026).

\bibitem{Keller2017}
D.~Keller, ``Modeling alignment enhancement for solid polarized targets,''
Eur.\ Phys.\ J.\ A \textbf{53}, 155 (2017).

\end{thebibliography}
\end{document}
